\documentclass[11pt,a4paper]{article}
\usepackage[UTF8]{ctex}
\usepackage[margin=2.5cm]{geometry}
\usepackage{booktabs}
\usepackage{longtable}
\usepackage{listings}
\usepackage{xcolor}
\usepackage{hyperref}
\usepackage{enumitem}
\usepackage{amssymb}
\hypersetup{colorlinks=true,linkcolor=blue,urlcolor=blue}

\lstset{
  basicstyle=\ttfamily\small,
  breaklines=true,
  frame=single,
  backgroundcolor=\color{gray!10},
  keywordstyle=\color{blue},
  commentstyle=\color{gray},
  language=bash
}

\title{\textbf{LLM Context Compaction Proxy}\\[0.3em]
\large 永不超限,永不停歇 · Never Full, Never Stop\\[0.3em]
\large 中文使用指南}
\author{LLM Context Compaction Proxy Contributors}
\date{2026 年 8 月}

\begin{document}

\maketitle
\tableofcontents
\newpage

% ============================================================
\section{功能亮点}
% ============================================================

\subsection{上下文压缩引擎}

\begin{longtable}{@{}lp{10cm}@{}}
  \toprule
  功能 & 说明 \\
  \midrule
  预压缩 & 上下文达 80\% 阈值时提前触发,避免触墙 \\
  溢出恢复 & 溢出错误时自动压缩并重试 \\
  增量压缩 & 只摘要自上次压缩后的新增消息(CoMem) \\
  并行压缩 & 大上下文分块并行压缩 \\
  并行摘要合并 & 将并行分块摘要合并为一份连贯摘要 \\
  双层压缩 & 网关层(85\% 缩减)+ Agent 层(L1 的 50\%),三种降级路径 \\
  CJK 感知 Token 估算 & 精确统计中文/英文/其他字符 token 数(两阶段) \\
  压缩安全验证 & 压缩结果必须小于原文,否则降级为激进截断 \\
  智能截断 & 基于角色的预算分配(头尾拆分) \\
  标识符保留 & 压缩时保持代码标识符完整 \\
  强制标识符列表 & 摘要中原样保留 UUID/哈希/URL/文件名 \\
  压缩缓存 & 30 分钟 TTL,指令感知 salt 键 \\
  压缩项剪枝 & 清理上次压缩遗留的低价值产物 \\
  查询位置优化 & 查询置尾、工具对邻接、近期窗口重排 \\
  缓存优化排序 & 层哈希感知重排 + \texttt{cache\_control} 断点 \\
  结构感知工具输出压缩 & 代码/日志/JSON 专用压缩器 \\
  手动压缩(\texttt{/compact}) & 类似 Claude Code 的 /compact,可强制压缩任意会话 \\
  激进截断兜底 & 压缩失败或膨胀时降级为截断 \\
  \bottomrule
\end{longtable}

\subsection{记忆与会话}

\begin{longtable}{@{}lp{10cm}@{}}
  \toprule
  功能 & 说明 \\
  \midrule
  语义记忆 & 情节—语义双层记忆(arXiv:2605.17625),6 类知识槽位 \\
  跨会话记忆 + FTS5 & SQLite 会话持久化 + 全文搜索 \\
  用户画像记忆 & 持久化用户偏好(USER.md),有界持久化 \\
  会话转录存档 & 保存/恢复原始转录以支持会话恢复 \\
  会话恢复 & 从转录 + 摘要恢复会话(\texttt{POST /sessions/\{id\}/resume}) \\
  背景知识注入 & 将近期会话知识注入压缩提示 \\
  会话上下文注入 & 自动向出站请求注入会话上下文 \\
  ARC 地址化引用 & 长工具结果替换为 ID 引用 \\
  ARC 持久化回查 & ARC 条目持久化到数据库,经 \texttt{GET /arc/\{arc\_id\}} 回查 \\
  承诺提取(CCL) & 提取并保留目标/约束/决策/错误 \\
  LLM 记忆提取 & LLM 驱动知识提取,正则兜底 \\
  思考掩码 & 剥离推理内容以节省 token \\
  密钥脱敏 & 自动脱敏 API Key/JWT/密码/Token \\
  孤儿工具对清理 & 清理悬空的 tool\_call/tool\_result 对 \\
  工具对邻接修复 & 修复非邻接的工具调用/结果顺序 \\
  \bottomrule
\end{longtable}

\subsection{Provider 与协议}

\begin{longtable}{@{}lp{10cm}@{}}
  \toprule
  功能 & 说明 \\
  \midrule
  Provider 抽象层 & 通用 OpenAI/Anthropic/Gemini 兼容(ABC) \\
  自动 Provider 探测 & 依据模型名/请求头/URL 自动识别格式 \\
  OpenAI$\rightarrow$Anthropic 请求转换 & 完整字段映射,含 tool\_calls $\rightarrow$ tool\_use \\
  Anthropic$\rightarrow$OpenAI 响应转换 & content/thinking/tool\_use 映射回 OpenAI \\
  SSE 流式转换 & Anthropic SSE $\rightarrow$ OpenAI SSE 逐事件转换 \\
  Anthropic 仅思考重试 & 响应仅含空思考块时自动流式重试 \\
  独立压缩 Provider & 压缩模型可使用独立上游与密钥 \\
  Gemini 密钥走请求头 & \texttt{x-goog-api-key} header,绝不进 URL/日志 \\
  模型注册表 & 67+ 已知上下文窗口的模型 \\
  透传路由 & 任意未匹配路径透明转发上游 \\
  \bottomrule
\end{longtable}

\subsection{可靠性与安全}

\begin{longtable}{@{}lp{10cm}@{}}
  \toprule
  功能 & 说明 \\
  \midrule
  熔断器 & 3 次失败 $\rightarrow$ 60 秒冷却,三状态机 \\
  抖动检测 & 检测压缩循环(3 次 / 5 条消息)$\rightarrow$ 激进截断 \\
  端点认证 & 22+ 端点受保护;无密钥时仅回环放行 \\
  并发安全 & 语义记忆线程锁、无竞态提示词(参数传递) \\
  压缩前后 Hooks & 压缩前后的 HTTP webhook(可阻断) \\
  健康与指标端点 & \texttt{/health}、\texttt{/metrics}(Prometheus 风格计数器) \\
  自适应保留轮次 & 对话过短时自动下调保留轮次 \\
  收缩窗口重试 & 每次重试递减保留轮次 \\
  \bottomrule
\end{longtable}

\subsection{MemSkill(自演进技能,V7)}

\begin{longtable}{@{}lp{10cm}@{}}
  \toprule
  功能 & 说明 \\
  \midrule
  技能注册表 & CRUD + 激活生命周期 + 快照回滚(持久化到 SQLite) \\
  技能控制器 & 关键词匹配 + Gumbel-Top-K 选择 \\
  参数覆盖管线 & 每技能的重要性权重 + 保真乘数 \\
  DELETE 型技能管线跳过 & 跳过 DELETE 操作的语义提取 \\
  技能性能追踪 & 每技能奖励/成功统计(\texttt{/skills/\{id\}/performance}) \\
  技能轨迹 & 近期压缩轨迹(\texttt{/skills/trajectories}) \\
  技能设计器 & 自动设计新技能(\texttt{/skills/designer/trigger}) \\
  \bottomrule
\end{longtable}

\subsection{HTTP API 端点(31 个)}

\begin{longtable}{@{}llp{7.5cm}@{}}
  \toprule
  方法 & 端点 & 说明 \\
  \midrule
  POST & \texttt{/chat/completions}(+\texttt{/v1/chat/completions}) & OpenAI 兼容对话(流式与非流式) \\
  POST & \texttt{/v1/messages} & Anthropic Messages API \\
  POST & \texttt{/compact} & 手动压缩会话 \\
  POST & \texttt{/summarize} & 供 CompactionProvider 插件使用的纯摘要端点 \\
  POST & \texttt{/session} & 创建会话 \\
  GET & \texttt{/sessions/search} & FTS5 全文会话搜索 \\
  GET & \texttt{/sessions/recent} & 近期会话 \\
  GET & \texttt{/sessions/\{id\}} & 会话详情 \\
  GET & \texttt{/sessions/\{id\}/transcript} & 原始转录导出 \\
  POST & \texttt{/sessions/\{id\}/resume} & 恢复会话 \\
  GET/POST & \texttt{/profile} & 用户画像读取/设置 \\
  GET/DELETE & \texttt{/memory} & 语义记忆读取/清空 \\
  GET & \texttt{/arc/\{arc\_id\}} & ARC 引用回查 \\
  GET & \texttt{/skills}、\texttt{/skills/\{id\}} & 技能列表/详情 \\
  POST & \texttt{/skills} & 创建技能(草稿) \\
  POST & \texttt{/skills/\{id\}/activate}、\texttt{/deprecate} & 技能生命周期 \\
  POST & \texttt{/skills/\{id\}/rollback} & 回滚到快照版本 \\
  GET & \texttt{/skills/\{id\}/performance} & 技能奖励统计 \\
  GET & \texttt{/skills/trajectories} & 压缩轨迹 \\
  POST & \texttt{/skills/designer/trigger} & 触发技能设计器 \\
  GET & \texttt{/health}、\texttt{/metrics} & 健康与指标 \\
  ANY & \texttt{/\{path:.*\}} & 透传到上游 \\
  \bottomrule
\end{longtable}

% ============================================================
\section{支持的 Agent(OpenClaw / Claude Code / Hermes / DeepSeek Harness)}
% ============================================================

代理已实测兼容 \textbf{OpenClaw、Claude Code、Hermes 与 DeepSeek Harness}
四大 Agent 框架。接入后上下文压缩完全由代理接管,Agent 自身不再压缩,
可无限期运行。

\subsection{OpenClaw}

在 \texttt{openclaw.json} 中将模型 provider 指向代理,并启用
\texttt{v6-compaction-provider} 插件:

\begin{lstlisting}[language={}]
{
  "models": {
    "providers": {
      "compaction-proxy": {
        "baseUrl": "http://127.0.0.1:8198",
        "api": "openai-completions",
        "models": [{ "id": "xopdeepseekv4flash" }]
      }
    }
  },
  "plugins": {
    "entries": {
      "v6-compaction-provider": {
        "enabled": true,
        "config": { "proxyUrl": "http://127.0.0.1:8198", "model": "xsparkx2agent" }
      }
    }
  }
}
\end{lstlisting}

已验证:OpenClaw 的 \texttt{xunfei} provider 已指向 \texttt{127.0.0.1:8198},
\texttt{v6-compaction-provider} 插件经 \texttt{/summarize} 走代理;
OpenClaw 原生 compaction 已关闭(\texttt{compaction.enabled=false}),
压缩完全由代理接管。

\subsection{Claude Code(Anthropic)}

\begin{lstlisting}
export ANTHROPIC_BASE_URL=http://localhost:8198
claude
\end{lstlisting}

或在 \texttt{\~/.claude/settings.json} 中:

\begin{lstlisting}[language={}]
{
  "env": { "ANTHROPIC_BASE_URL": "http://localhost:8198" },
  "autoCompactEnabled": false
}
\end{lstlisting}

已验证:\texttt{POST /v1/messages} 经代理路由(缺 key 返回 401 证明
Anthropic 格式识别成功);\texttt{autoCompactEnabled: false} 关闭
Claude Code 自带压缩。

\subsection{Hermes Agent}

编辑 \texttt{\~/.hermes/config.yaml}:

\begin{lstlisting}[language={}]
model:
  api_mode: chat_completions          # OpenAI 格式
  base_url: http://127.0.0.1:8198/v1  # 指向代理
  default: xopglm51                   # 上游可用模型均可
\end{lstlisting}

已验证:Hermes 的 \texttt{chat\_completions} 模式与代理
\texttt{/v1/chat/completions} 完全兼容;Hermes 自带压缩已关闭
(\texttt{compression.enabled: false})。

\subsection{DeepSeek Harness}

编辑 \texttt{\$DSH\_HOME/settings.yaml}(默认 \texttt{\~/.dsh/settings.yaml}):

\begin{lstlisting}[language={}]
llm-pi-ai:
  providers:
    xfyun-maas:
      apiKeyEnv: XF_MASS_API_KEY
      api: openai-completions           # OpenAI 兼容协议
      baseURL: http://127.0.0.1:8198    # 指向代理
      models:
        - id: xsparkx2agent
          contextWindow: 131072
          maxTokens: 8192
\end{lstlisting}

已验证:Harness 的 \texttt{xfyun-maas} provider 已指向
\texttt{127.0.0.1:8198};Harness 原生压缩经 \texttt{cordis.patch.yml}
关闭(\texttt{dsh-compaction-basic} $\rightarrow$ \texttt{auto: false})。

% ============================================================
\section{安装和使用说明}
% ============================================================

\subsection{前置条件}

\begin{itemize}
  \item Python 3.10+
  \item \texttt{aiohttp}(\texttt{pip install aiohttp})
\end{itemize}

\subsection{安装步骤}

\textbf{1. 克隆仓库}

\begin{lstlisting}
git clone https://github.com/061115xhsm/NeverFull-NeverStop-LLM-Context-Compaction-Proxy.git
cd NeverFull-NeverStop-LLM-Context-Compaction-Proxy
\end{lstlisting}

\textbf{2. 安装依赖}

\begin{lstlisting}
pip install aiohttp
\end{lstlisting}

\textbf{3. 配置上游并运行}

\begin{lstlisting}
export COMPACTION_PROXY_UPSTREAM=https://api.openai.com/v1
export COMPACTION_PROXY_MODEL=gpt-4o-mini
python3 compaction-proxy.py
\end{lstlisting}

代理监听在 \texttt{localhost:8198}。将 Agent 的 API 端点指向此处即可。

\subsection{接入你的 LLM 提供商}

\begin{longtable}{@{}lp{9cm}@{}}
  \toprule
  提供商 & 配置 \\
  \midrule
  OpenAI / 兼容系 & \texttt{export COMPACTION\_PROXY\_UPSTREAM=https://api.openai.com/v1} \\
  Anthropic & \texttt{.../api.anthropic.com} + \texttt{UPSTREAM\_IS\_ANTHROPIC=true} \\
  Gemini & \texttt{.../generativelanguage.googleapis.com} \\
  DeepSeek & \texttt{.../api.deepseek.com/v1} \\
  Ollama(本地) & \texttt{http://localhost:11434/v1}(默认) \\
  \bottomrule
\end{longtable}

\subsection{帮助与支持}

\begin{itemize}
  \item \textbf{问题追踪}:GitHub Issues
  \item \textbf{健康检查}:\texttt{curl http://localhost:8198/health}
  \item \textbf{文档}:仓库内的 \texttt{README.tex} 与 \texttt{docs/technical-document.tex}
\end{itemize}

% ============================================================
\section{示例和代码片段}
% ============================================================

\subsection{Python SDK 示例}

\begin{lstlisting}[language=Python]
import openai

client = openai.OpenAI(
    base_url="http://localhost:8198/v1",
    api_key="your-key"
)

response = client.chat.completions.create(
    model="gpt-4o-mini",
    messages=[{"role": "user", "content": "Hello!"}],
    stream=True
)
\end{lstlisting}

创建指向 \texttt{localhost:8198/v1} 的 OpenAI 客户端,调用
\texttt{chat.completions.create()} 发送消息并以流式方式读取响应。

\subsection{验证代理是否生效}

\begin{lstlisting}
# Health check
curl http://localhost:8198/health

# Send a test request through the proxy
curl -X POST http://localhost:8198/v1/chat/completions \
  -H "Content-Type: application/json" \
  -H "Authorization: Bearer YOUR_KEY" \
  -d '{"model":"gpt-4o-mini","messages":[{"role":"user","content":"Hello"}]}'
\end{lstlisting}

% ============================================================
\section{项目结构和文件组织}
% ============================================================

\begin{longtable}{@{}llp{6.5cm}@{}}
  \toprule
  文件/目录 & 描述 & 用途 \\
  \midrule
  \texttt{compaction-proxy.py} & 主程序 & 代理的全部核心逻辑 \\
  \texttt{README.md} / \texttt{README.tex} & 项目说明 & 使用指南与文档 \\
  \texttt{README-zh.md} / \texttt{README-zh.tex} & 中文说明 & 中文使用指南 \\
  \texttt{docs/} & 文档目录 & 技术文档与接入指南 \\
  \texttt{.env.example} & 配置模板 & 环境变量配置示例 \\
  \texttt{systemd/} & 服务模板 & systemd 用户服务单元 \\
  \texttt{LICENSE} & 许可证 & MIT 许可证文本 \\
  \bottomrule
\end{longtable}

\subsection{主程序(\texttt{compaction-proxy.py})}

包含代理的全部核心逻辑:Provider 抽象层(OpenAI/Anthropic/Gemini)、压缩引擎、
语义记忆、会话存储、技能注册表、熔断器与抖动检测等。代码按模块化组织,
每个类与函数职责单一。

\subsection{文档目录(\texttt{docs/})}

包含技术文档与接入指南,帮助用户和开发者理解与使用项目。

\subsection{测试与配置}

\texttt{.env.example} 提供完整配置模板,\texttt{systemd/compaction-proxy.service}
提供系统服务化部署模板,方便生产环境使用。

% ============================================================
\section{贡献指南}
% ============================================================

贡献开源项目不仅包括代码,还包括文档更新、问题报告、新功能建议等。

\subsection{贡献流程}

\begin{enumerate}[leftmargin=2em]
  \item \textbf{了解项目}:阅读文档和代码,了解项目的目标、架构和设计原则;
  \item \textbf{找到贡献机会}:查看问题跟踪器,找到可解决的问题;关注未来计划与里程碑;
  \item \textbf{贡献代码}:Fork 项目 $\rightarrow$ 本地开发测试 $\rightarrow$ 提交 Pull Request。
\end{enumerate}

\subsection{提交代码的标准}

\begin{itemize}
  \item 代码必须符合项目的编码标准和风格指南;
  \item 提交的代码必须通过所有测试;
  \item 代码应包含单元测试,确保功能的正确性。
\end{itemize}

\subsection{提交问题与拉取请求}

\begin{longtable}{@{}lp{5.5cm}p{6cm}@{}}
  \toprule
  方面 & 提交问题 & 提交拉取请求 \\
  \midrule
  标题 & 明确、具体 & 清晰、描述目的 \\
  描述 & 详细、包含重现步骤 & 详细、解释更改的必要性 \\
  附加信息 & 屏幕截图、动画 & 符合编码和风格标准 \\
  \bottomrule
\end{longtable}

% ============================================================
\section{许可证}
% ============================================================

本项目采用 \textbf{MIT License},详见仓库根目录的 \texttt{LICENSE} 文件。

MIT 许可证允许他人自由使用、修改和分发你的代码,只要他们包含原始许可证
和版权声明。

\subsection{许可证的类型}

\begin{itemize}
  \item \textbf{MIT License}:宽松许可,允许自由使用/修改/分发,保留版权声明即可;
  \item \textbf{GNU GPL}:要求使用、修改或分发该许可证下代码的人都必须将其更改公开,
        并使用相同许可证。
\end{itemize}

% ============================================================
\section{联系信息和致谢}
% ============================================================

\subsection{联系信息}

如果你有任何问题或建议,请随时通过 GitHub Issue 与我们联系,或在项目仓库
的 Discussions 中发起讨论。

\subsection{致谢}

感谢所有为本项目贡献过代码、文档、测试与建议的开发者与用户。你的每一次
提交、每一个 Issue、每一条评论,都在推动这个项目向前。

\subsection{社区文化}

我们欢迎每一个人的参与和贡献,无论你的技能水平如何,都有你的一席之地。

% ============================================================
\section*{结语}
% ============================================================

LLM Context Compaction Proxy 是一段探索之旅的产物:从研究论文中的压缩技术,
到生产环境中的工程实践;从解决自己的上下文超限之痛,到帮助更多开发者的
Agent 无限期运行。愿这份中文指南能成为你的灯塔,指引你在上下文压缩的世界里,
找到属于自己的路径。

\vspace{1em}
\noindent \emph{—— LLM Context Compaction Proxy Contributors}

\end{document}
